% =============================================================================
\section{Problem Formulation and Assumptions}\label{sec:formulation}
% =============================================================================

We seek to formalize and estimate an \emph{effective context length} for an embedding-generating model $f$, relative to (i) a reference position $r$, (ii) a distribution of sequences $P_X$, (iii) a perturbation mechanism $\Pi$, and (iv) an embedding discrepancy $d_\cZ$. Formally, we estimate a functional
\[
\Psi(f, P_X, \Pi, d_\cZ, \beta) \;=\; \ECL_\beta(r).
\]

\subsection{Core question}

Intuitively, ECL should be small if embeddings at $r$ depend mostly on nearby positions, and large if distal context materially affects embeddings. We pose the following estimation problem:

\begin{quote}
\emph{Given black-box access to $f$ and the ability to sample sequences $X\sim P_X$ and perturbations $X^{(S)}\sim \Pi_S(\cdot\mid X)$, estimate the minimal context radius $\ell$ such that perturbations outside $W_\ell(r)$ contribute at most a prescribed fraction of total embedding variability.}
\end{quote}

\subsection{Assumptions}

We state explicit conditions under which our theoretical results hold. Many can be weakened; clarity is prioritized.

\begin{assumption}[Measurability and finite second moment]\label{ass:meas}
$f$ is measurable, $d_\cZ$ is measurable, and
$\E\big[d_\cZ(f(X),f(X'))\big] < \infty$
for $X,X'$ distributed as specified by the perturbation mechanism.
\end{assumption}

\begin{assumption}[Perturbation locality]\label{ass:local}
For each set $S$, $\Pi_S(\cdot\mid x)$ modifies only coordinates in $S$: $X^{(S)}_{-S}=x_{-S}$ almost surely.
\end{assumption}

\begin{assumption}[Conditional local independence]\label{ass:condind}
When needed for Sobol-type decompositions (specifically, \cref{thm:sobol_equiv}), we assume that conditioned on a local context window of size $k$ (e.g., $k=2$ for dinucleotide context), coordinates are approximately independent:
\[
P(X_i \mid X_{-i}) \approx P(X_i \mid X_{i-k:i+k}).
\]
This is biologically more realistic than full coordinate independence and is satisfied by $k$-th order Markov models for DNA.
\end{assumption}

\begin{assumption}[Bounded influence variables]\label{ass:subgauss}
For each position $i$, the random variable $U_i := d_\cZ(f(X), f(X^{(i)}))$ satisfies $0 \le U_i \le M$ almost surely for some constant $M < \infty$, with finite variance $\Var(U_i) = \sigma_i^2 < \infty$. Boundedness holds whenever the embedding map $f$ is Lipschitz (\cref{ass:lipschitz}) and inputs are drawn from a finite alphabet, which is the standard setting in genomics.
\end{assumption}

\begin{assumption}[Lipschitz stability]\label{ass:lipschitz}
There exists $K>0$ such that for all $x,x'\in\cX$,
$\norm{f(x)-f(x')}_\cZ \le K \cdot \Ham(x,x')$,
where $\Ham$ is Hamming distance. Under this assumption with bounded inputs, $U_i = d_\cZ(f(X),f(X^{(i)})) \le K^2$ almost surely under squared Euclidean discrepancy, supplying the boundedness constant $M$ used in \cref{ass:subgauss}.
\end{assumption}
